% ============================================================
% PARTE AGUSTÍN — ejemplos con parámetros, ejemplo con monopolio
% y conclusiones  (Secciones IV, V y VI del paper)
% ============================================================
% Las slides de esta parte son APOYO VISUAL del guion:
% "Nota - Guion presentación (Secciones IV, V y VI).md"
% Frases sueltas (sin bullets salvo enumeraciones); cada tabla con su pie
% estilo paper ("Nota. ..."); la lectura en vivo la hace el orador.

% Cabecera de tres filas para las Tablas I y II, como el paper:
% delta_v (grupos) / rho_vf (columnas) / #1 = variable de fila del bloque
\newcommand{\cabTablas}[1]{%
  \rowcolor{gristitulo}
  \multicolumn{1}{l}{$\delta_v$}
  & \multicolumn{3}{c}{$0.05$}
  & \multicolumn{3}{c}{$0.10$}
  & \multicolumn{3}{c}{$0.25$} \\
  \rowcolor{gristitulo}
  \multicolumn{1}{l}{$\rho_{vf}$}
  & $-0.5$ & $0.0$ & $0.5$
  & $-0.5$ & $0.0$ & $0.5$
  & $-0.5$ & $0.0$ & $0.5$ \\
  \rowcolor{gristitulo}
  \multicolumn{1}{l}{#1} & & & & & & & & & \\}

% (Sección IV del paper — guion §1 y §2)
\begin{frame}{Aplicación del modelo con ejemplos numéricos}
  Todas las tablas se evalúan en el punto $\rj{V} = \az{F} = 1$: el primer punto
  donde la regla clásica del VAN acepta la inversión.
  \begin{ecmed}
    \az{\sigma_v^2 = \sigma_f^2 = 0.04}, \qquad
    \az{\delta_v = \delta_f = 0.10}, \qquad
    \az{\rho_{vf} = 0}, \qquad \az{\lambda = 0}
  \end{ecmed}
  Datos para el caso base:
  \begin{itemize}
    \item \az{$\sigma = 0.20$}: desvío del equity desapalancado de EEUU.
    \item \az{$\delta_v = 0.10$}: ratio ganancias/precio de un proyecto
          instalado.
    \item \az{$\delta_f = 0.10$}: tasa libre de riesgo, con \az{$F$}
          determinístico.
  \end{itemize}
\end{frame}

% Tabla I — bloque 1 (varianza)
\begin{frame}{Valor de la oportunidad de inversión}
  \vspace{2mm}
  {\centering
  {\fontsize{8}{11}\selectfont
   \renewcommand{\arraystretch}{1.3}%
   \setlength{\tabcolsep}{4pt}%
   \arrayrulecolor{grispie}%
   \begin{tabular}{>{\bfseries\color{navy}}l ccccccccc}
     \cabTablas{$\sigma_v^2,\,\sigma_f^2$}
     0.01 & 0.33 & 0.30 & 0.28 & 0.14 & 0.12 & 0.08 & 0.03 & 0.02 & 0.01 \\ \hline
     0.02 & 0.38 & 0.34 & 0.30 & 0.20 & 0.16 & 0.12 & 0.06 & 0.04 & 0.02 \\ \hline
     0.04 & \textbf{0.45} & \textbf{0.40} & \textbf{0.34} & \textbf{0.27} &
     \cellcolor{celeste}\textbf{0.23} & \textbf{0.16} & \textbf{0.11} &
     \textbf{0.08} & \textbf{0.04} \\ \hline
     0.10 & 0.57 & 0.51 & 0.43 & 0.40 & 0.34 & 0.25 & 0.21 & 0.16 & 0.09 \\ \hline
     0.20 & 0.67 & 0.61 & 0.51 & 0.52 & 0.45 & 0.34 & 0.32 & 0.25 & 0.16 \\ \hline
     0.30 & 0.73 & 0.67 & 0.57 & 0.60 & 0.52 & 0.40 & 0.39 & 0.32 & 0.21 \\
   \end{tabular}}\par}
  \pietabla{Valor de la oportunidad de inversión cuando $V = F = 1$, por
    dólar de proyecto, calculado con las ecuaciones (4) y (12). Bloque 1:
    varianza ($\lambda = 0$, $\delta_f = 0.10$ fijos). Parámetros del caso
    base: $\sigma_v^2 = \sigma_f^2 = 0.04$; $\delta_v = \delta_f = 0.10$;
    $\lambda = 0$.}
\end{frame}

% Tabla I — bloques 2 y 3, cada uno con su pie
\begin{frame}{Valor de la oportunidad de inversión}
  \vspace{1mm}
  {\centering
  {\fontsize{8}{11}\selectfont
   \renewcommand{\arraystretch}{1.08}%
   \setlength{\tabcolsep}{4pt}%
   \arrayrulecolor{grispie}%
   \begin{tabular}{>{\bfseries\color{navy}}l ccccccccc}
     \cabTablas{$\lambda$}
     0.00 & \textbf{0.45} & \textbf{0.40} & \textbf{0.34} & \textbf{0.27} &
     \cellcolor{celeste}\textbf{0.23} & \textbf{0.16} & \textbf{0.11} &
     \textbf{0.08} & \textbf{0.04} \\ \hline
     0.05 & 0.33 & 0.29 & 0.24 & 0.23 & 0.19 & 0.13 & 0.10 & 0.07 & 0.04 \\ \hline
     0.10 & 0.27 & 0.23 & 0.19 & 0.20 & 0.16 & 0.12 & 0.10 & 0.07 & 0.04 \\ \hline
     0.25 & 0.19 & 0.16 & 0.12 & 0.15 & 0.12 & 0.09 & 0.09 & 0.06 & 0.04 \\
   \end{tabular}}\par}
  \pietabla{Bloque 2: probabilidad anual $\lambda$ de que $V$ salte a cero
    ($\sigma_v^2 = \sigma_f^2 = 0.04$, $\delta_f = 0.10$ fijos).}
  \vspace{1.5mm}
  {\centering
  {\fontsize{8}{11}\selectfont
   \renewcommand{\arraystretch}{1.08}%
   \setlength{\tabcolsep}{4pt}%
   \arrayrulecolor{grispie}%
   \begin{tabular}{>{\bfseries\color{navy}}l ccccccccc}
     \cabTablas{$\delta_f$}
     0.01 & 0.30 & 0.23 & 0.14 & 0.18 & 0.13 & 0.07 & 0.08 & 0.06 & 0.03 \\ \hline
     0.05 & 0.37 & 0.31 & 0.23 & 0.22 & 0.17 & 0.10 & 0.09 & 0.06 & 0.03 \\ \hline
     0.10 & \textbf{0.45} & \textbf{0.40} & \textbf{0.34} & \textbf{0.27} &
     \cellcolor{celeste}\textbf{0.23} & \textbf{0.16} & \textbf{0.11} &
     \textbf{0.08} & \textbf{0.04} \\ \hline
     0.25 & 0.60 & 0.58 & 0.56 & 0.42 & 0.39 & 0.36 & 0.18 & 0.15 & 0.10 \\
   \end{tabular}}\par}
  \pietabla{Bloque 3: $\delta_f$ ($\sigma_v^2 = \sigma_f^2 = 0.04$,
    $\lambda = 0$ fijos).}
\end{frame}

% Tabla II — bloque 1 (varianza)
\begin{frame}{Nivel de $V/F$ al que conviene invertir}
  \vspace{1mm}
  {\centering
  {\fontsize{8}{11}\selectfont
   \renewcommand{\arraystretch}{1.25}%
   \setlength{\tabcolsep}{4pt}%
   \arrayrulecolor{grispie}%
   \begin{tabular}{>{\bfseries\color{navy}}l ccccccccc}
     \cabTablas{$\sigma_v^2,\,\sigma_f^2$}
     0.01 & 2.50 & 2.35 & 2.18 & 1.47 & 1.37 & 1.25 & 1.09 & 1.06 & 1.03 \\ \hline
     0.02 & 2.91 & 2.64 & 2.35 & 1.72 & 1.56 & 1.37 & 1.18 & 1.12 & 1.06 \\ \hline
     0.04 & \textbf{3.65} & \textbf{3.17} & \textbf{2.64} & \textbf{2.13} &
     \cellcolor{celeste}\textbf{1.86} & \textbf{1.56} & \textbf{1.34} &
     \textbf{1.24} & \textbf{1.12} \\ \hline
     0.10 & 5.65 & 4.56 & 3.41 & 3.19 & 2.62 & 2.00 & 1.77 & 1.54 & 1.29 \\ \hline
     0.20 & 8.77 & 6.70 & 4.56 & 4.79 & 3.73 & 2.62 & 2.44 & 2.00 & 1.54 \\ \hline
     0.30 & 11.83 & 8.77 & 5.65 & 6.34 & 4.79 & 3.19 & 3.07 & 2.44 & 1.77 \\
   \end{tabular}}\par}
  \pietabla{Valor de los beneficios relativo al costo de inversión ($V/F$)
    al que la inversión es óptima ($C^*$), calculado con las ecuaciones (4)
    y (12). Bloque 1: varianza ($\lambda = 0$, $\delta_f = 0.10$ fijos).}
  \begin{cajatexto}
    Con parámetros razonables es óptimo esperar hasta que el valor del
    proyecto \textbf{duplique} el costo de inversión.
  \end{cajatexto}
\end{frame}

% Tabla II — bloques 2 y 3, cada uno con su pie
\begin{frame}{Nivel de $V/F$ al que conviene invertir}
  \vspace{1mm}
  {\centering
  {\fontsize{8}{11}\selectfont
   \renewcommand{\arraystretch}{1.08}%
   \setlength{\tabcolsep}{4pt}%
   \arrayrulecolor{grispie}%
   \begin{tabular}{>{\bfseries\color{navy}}l ccccccccc}
     \cabTablas{$\lambda$}
     0.00 & \textbf{3.65} & \textbf{3.17} & \textbf{2.64} & \textbf{2.13} &
     \cellcolor{celeste}\textbf{1.86} & \textbf{1.56} & \textbf{1.34} &
     \textbf{1.24} & \textbf{1.12} \\ \hline
     0.05 & 2.50 & 2.23 & 1.92 & 1.86 & 1.67 & 1.44 & 1.32 & 1.23 & 1.12 \\ \hline
     0.10 & 2.10 & 1.90 & 1.67 & 1.72 & 1.56 & 1.37 & 1.30 & 1.22 & 1.12 \\ \hline
     0.25 & 1.67 & 1.54 & 1.40 & 1.51 & 1.40 & 1.27 & 1.27 & 1.19 & 1.11 \\
   \end{tabular}}\par}
  \pietabla{Bloque 2: probabilidad anual $\lambda$ de que $V$ salte a cero
    ($\sigma_v^2 = \sigma_f^2 = 0.04$, $\delta_f = 0.10$ fijos).}
  \vspace{1.5mm}
  {\centering
  {\fontsize{8}{11}\selectfont
   \renewcommand{\arraystretch}{1.08}%
   \setlength{\tabcolsep}{4pt}%
   \arrayrulecolor{grispie}%
   \begin{tabular}{>{\bfseries\color{navy}}l ccccccccc}
     \cabTablas{$\delta_f$}
     0.01 & 2.31 & 1.89 & 1.46 & 1.64 & 1.43 & 1.22 & 1.25 & 1.17 & 1.08 \\ \hline
     0.05 & 2.85 & 2.38 & 1.86 & 1.83 & 1.58 & 1.32 & 1.28 & 1.19 & 1.10 \\ \hline
     0.10 & \textbf{3.65} & \textbf{3.17} & \textbf{2.64} & \textbf{2.13} &
     \cellcolor{celeste}\textbf{1.86} & \textbf{1.56} & \textbf{1.34} &
     \textbf{1.24} & \textbf{1.12} \\ \hline
     0.25 & 6.42 & 5.96 & 5.49 & 3.35 & 3.09 & 2.81 & 1.62 & 1.49 & 1.33 \\
   \end{tabular}}\par}
  \pietabla{Bloque 3: $\delta_f$ ($\sigma_v^2 = \sigma_f^2 = 0.04$,
    $\lambda = 0$ fijos).}
\end{frame}

% (guion §3 "qué señalar" + §4 — síntesis de patrones)
\begin{frame}{Estáticas comparativas de las Tablas I y II}
  \begin{itemize}
    \item \az{$\sigma^2$} $\uparrow$\;: valor de la opción $\uparrow$,
          \rj{$C^*$} $\uparrow$.
    \item \az{$\delta_f$} $\uparrow$\;: valor $\uparrow$, \rj{$C^*$}
          $\uparrow$ --- crece el ahorro por diferir el desembolso.
    \item \rj{$\delta_v$} $\uparrow$\;: valor $\downarrow$, \rj{$C^*$}
          $\downarrow$ --- el efecto más fuerte.
    \item \az{$\rho_{vf}$} $\uparrow$\;: valor $\downarrow$, \rj{$C^*$}
          $\downarrow$ --- se reduce la varianza del ratio.
    \item \az{$\lambda$} $\uparrow$\;: valor $\downarrow$, \rj{$C^*$}
          $\downarrow$ --- equivale a descontar a \az{$\mu + \lambda$}.
  \end{itemize}
\end{frame}

% Tabla III — OPCIONAL: sacar si la parte queda larga (guion §5)
\begin{frame}{El peso de la incertidumbre en el valor}
  \vspace{2mm}
  {\centering
  {\fontsize{8.5}{11.5}\selectfont
   \renewcommand{\arraystretch}{1.3}%
   \setlength{\tabcolsep}{8pt}%
   \arrayrulecolor{grispie}%
   \begin{tabular}{>{\bfseries\color{navy}}l cccc}
     \rowcolor{gristitulo}
     \multicolumn{1}{l}{$\sigma^2$} & $0.02$ & $0.04$ & $0.10$ & $0.30$ \\
     \rowcolor{gristitulo}
     \multicolumn{1}{l}{$\delta_v$} & & & & \\
     0.02 & 5.3\%  & 9.4\%  & 17.2\% & 28.4\% \\ \hline
     0.04 & 12.5\% & 20.1\% & 32.6\% & 47.1\% \\ \hline
     0.06 & 25.6\% & 36.4\% & 50.8\% & 64.6\% \\ \hline
     0.08 & 51.4\% & 61.9\% & 73.0\% & 82.1\% \\ \hline
     0.10 & 100\%  & 100\%  & 100\%  & 100\%  \\
   \end{tabular}}\par}
  \pietabla{Porcentaje del valor de la opción atribuible a la incertidumbre,
    comparando la ecuación (4) con la (13) determinística ($V = F = 1$,
    $\delta_f = 0.10$).}
  \vspace{2.5mm}
  El cálculo mantiene fijos los retornos requeridos al variar $\sigma^2$:
  válido solo si el riesgo adicional es \concepto{no sistemático}.
\end{frame}

% (Sección V del paper — guion §1 y §2)
\begin{frame}{El monopolista: tecnología, demanda y shock}
  El valor del proyecto se construye desde una firma explícita:
  \begin{ecmed}
    \rj{Q_t} = \az{K^{1-\beta} L_t^{\beta}}, \qquad
    \rj{P_t} = \az{\theta_t\, Q_t^{-1/\eta}}, \qquad
    \rj{\frac{d\theta}{\theta}} = \az{\alpha_\theta\,dt + \sigma_\theta\,dz_\theta}
  \end{ecmed}
  Capital rígido (la inversión irreversible), trabajo flexible; toda la
  incertidumbre entra por el desplazador de demanda \az{$\theta_t$}.
  \begin{ecmed}
    \rj{\pi_t^{*}} = \az{B\,\theta_t^{\gamma}}, \qquad
    \az{\gamma = \bigl[1 - \beta(1 - 1/\eta)\bigr]^{-1} > 1}
  \end{ecmed}
  El beneficio maximizado es convexo en el shock: con \az{$\beta = 2/3$} y
  \az{$\eta = 2$}, \az{$\gamma = 1.5$}.
\end{frame}

% (guion §3 — el momento estrella: delta_v derivado)
\begin{frame}{El monopolista: $V$ y $\delta_v$ endógenos}
  Valor del proyecto instalado y retornos:
  \begin{ecmed}
    \rj{V(\theta_0)} = \az{\frac{B\,\theta_0^{\gamma}}{\hat{\alpha}_v - \alpha_v}},
    \qquad
    \az{\alpha_v = \gamma\alpha_\theta + \tfrac{1}{2}\gamma(\gamma-1)\sigma_\theta^2},
    \qquad
    \az{\hat{\alpha}_v = r + \phi\,\rho_{\theta m}\,\gamma\,\sigma_\theta}
  \end{ecmed}
  \cajaec{\rj{\delta_v} = \az{\hat{\alpha}_v - \alpha_v}
    = \az{B\,\theta_0^{\gamma} / V_0}}
  El payout del proyecto instalado deja de ser un supuesto: queda derivado
  de la tecnología y la demanda.
\end{frame}

% Tabla IV (guion §4)
\begin{frame}{Valores de la opción y reglas de inversión}
  \vspace{1mm}
  {\centering
  {\fontsize{8}{11}\selectfont
   \renewcommand{\arraystretch}{1.25}%
   \setlength{\tabcolsep}{4pt}%
   \arrayrulecolor{grispie}%
   \begin{tabular}{>{\bfseries\color{navy}}l ccccccccc}
     \rowcolor{gristitulo}
     \multicolumn{1}{l}{$\alpha_\theta$}
     & \multicolumn{3}{c}{$-0.03$}
     & \multicolumn{3}{c}{$0.00$}
     & \multicolumn{3}{c}{$0.01$} \\
     \rowcolor{gristitulo}
     \multicolumn{1}{l}{$\rho_{\theta m}$}
     & $0.0$ & $0.2$ & $0.5$
     & $0.0$ & $0.2$ & $0.5$
     & $0.0$ & $0.2$ & $0.5$ \\
     \rowcolor{gristitulo}
     \multicolumn{1}{l}{$\sigma_\theta^2$} & & & & & & & & & \\
     0.01 & 0.074 & 0.060 & 0.046 & 0.189 & 0.131 & 0.084 & 0.290 & 0.189 & 0.111 \\
          & (1.22) & (1.18) & (1.13) & (1.68) & (1.43) & (1.26) & (2.25) & (1.68) & (1.35) \\ \hline
     0.04 & 0.223 & 0.166 & 0.118 & 0.413 & 0.266 & 0.166 & 0.549 & 0.326 & 0.191 \\
          & (1.85) & (1.58) & (1.38) & (3.28) & (2.09) & (1.58) & (5.28) & (2.50) & (1.69) \\ \hline
     0.10 & 0.440 & 0.304 & 0.206 & 0.755 & 0.430 & 0.259 & * & 0.497 & 0.284 \\
          & (3.58) & (2.35) & (1.76) & (13.7) & (3.47) & (2.05) & * & (4.36) & (2.20) \\ \hline
     0.20 & 0.765 & 0.474 & 0.305 & * & 0.632 & 0.363 & * & 0.715 & 0.387 \\
          & (14.6) & (4.02) & (2.35) & * & (7.37) & (2.80) & * & (11.0) & (3.02) \\
   \end{tabular}}\par}
  \pietabla{Valor de la opción; entre paréntesis, nivel óptimo de inversión.
    Parámetros: $\eta = 2$; $r = 0.05$; $\beta = 2/3$; $\phi = 0.50$;
    $V = 1$; $F = 1$ no estocástico. *Para esos valores $\delta_v < 0$ y el
    valor de la opción no está definido.}
\end{frame}

% (guion §6 — autocrítica del paper)
\begin{frame}{Los límites del supuesto browniano}
  La demanda es una variable real: bajo MBG su incertidumbre crece
  linealmente con el horizonte, sin límite.
  \par\vspace{3mm}
  Un proceso con \concepto{reversión a la media} sería más realista, pero
  no admite solución analítica.
  \par\vspace{3mm}
  Un payout \rj{$\delta_v$} creciente en \az{$V$} acotaría el valor y
  devolvería un \rj{$C^*$} finito.
  \begin{cajatexto}
    El MBG es un buen supuesto para valores presentes; aplicado a
    fundamentos reales puede producir resultados implausibles.
  \end{cajatexto}
\end{frame}

% (Sección VI del paper)
\begin{frame}{Conclusiones}
  \vspace{1mm}
  La regla clásica del VAN invierte \emph{demasiado pronto}: bajo incertidumbre
  e irreversibilidad, VAN $=0$ no es señal de invertir, sino de que todavía
  conviene esperar.
  \par\vspace{3mm}
  No alcanza con $\rj{V} > \az{F}$: hay que esperar hasta
  $\rj{V}/\az{F} \geq \rj{C^*}$ ($\approx 1.86$ en el caso base) --- el proyecto
  debe valer \resultado{un 86\% más} que su costo antes de invertir.
  \par\vspace{3mm}
  \begin{cajatexto}
    ``El valor perdido por adoptar subóptimamente un proyecto de VAN cero está
    fácilmente entre el \textbf{10 y el 20 por ciento} (o más) del valor del
    proyecto.''
  \end{cajatexto}
  \vspace{1mm}
  Bajo incertidumbre, esperar deja de ser indecisión: es una \resultado{decisión
  con valor propio y medible}.
\end{frame}

% ============================================================
% LITERATURA POSTERIOR (pedido por Mariano) — va al final, tras Conclusiones
% Las limitaciones (MBG, inversión indivisible, irreversibilidad) se sacaron de
% "Conclusiones" para no duplicar; se desarrollan acá.
% ============================================================
\begin{frame}{Limitaciones del modelo base}
  El marco de 1986 es potente, pero deja afuera tres situaciones que la
  literatura posterior abordó:
  \begin{itemize}
    \item \textbf{El costo $F$ como MBG perpetuo}: tratar un bien físico como
          un browniano sin límite es irreal; a largo plazo los costos de
          capital tienden a \emph{revertir a la media}.
    \item \textbf{Inversión indivisible (``lumpy'')}: no admite proyectos
          complejos donde el capital se desembolsa \emph{por etapas}.
    \item \textbf{Sin flexibilidad intermedia}: no contempla encoger, expandir
          ni abandonar parcialmente la escala si el mercado cambia.
  \end{itemize}
\end{frame}

% Inversión en fases — Majd \& Pindyck (1987), Pindyck (1993)
\begin{frame}{Inversión en fases: trabajos posteriores}
  Cuando el capital se desembolsa por etapas, invertir hoy compra el
  \concepto{derecho a seguir}, no la obligación de terminar el proyecto:
  \begin{itemize}
    \item \textbf{Fase de estudio}: un desembolso chico y reversible (un
          estudio de factibilidad) revela información sin comprometer gran
          capital.
    \item \textbf{Opciones compuestas}: cada fase es una \emph{opción sobre
          otra opción} --- completar la Fase 1 habilita, pero no obliga, a la
          Fase 2.
    \item \textbf{Umbral más bajo}: la información temprana reduce el riesgo y
          \resultado{baja el umbral óptimo} \rj{$C^*$}.
  \end{itemize}
  \par\vspace{2mm}
  {\fontsize{8}{10}\selectfont Majd y Pindyck (1987), ``Time to Build, Option
   Value, and Investment Decisions''; Pindyck (1993), ``Investments of
   Uncertain Cost''. \emph{Journal of Financial Economics}.\par}
\end{frame}

% Evolución post-1986 (OPCIONAL: sumar más referencias; sacar si queda largo)
\begin{frame}{La evolución de la teoría (post-1986)}
  \begin{itemize}
    \item \textbf{Dixit y Pindyck (1994)}: el libro \emph{Investment Under
          Uncertainty} unifica la teoría y la vuelve central en las finanzas
          corporativas.
    \item \textbf{Histéresis (Pindyck 1988; Dixit 1989)}: entrada y salida de
          mercados bajo incertidumbre; explica el retraso de las firmas ante
          cambios de precios.
    \item \textbf{Interacción de opciones (Trigeorgis 1993)}: proyectos con
          varias opciones simultáneas e interdependientes ---esperar, expandir,
          abandonar.
  \end{itemize}
\end{frame}
